\begin{ZhChapter}
\chapter{結論與未來展望}

\section{實驗成果總結}
\text 本研究針對現有RESTful API自動化測試工具在語意理解能力不足、跨操作依賴推理能力有限、測試流程恢復能力缺乏，以及缺少工作流程層級決策機制等問題，提出一套結合多代理協作、工作流程恢復層與學習型決策探索之RESTful API自動化測試框架。研究架構整合大型語言模型（Large Language Model, LLM）、多代理系統（Multi-Agent System, MAS）、AutoGen、LangGraph、模型上下文協定（Model Context Protocol, MCP）、代理人對代理人協定（Agent-to-Agent Protocol, A2A）、離線強化學習（Offline Reinforcement Learning）以及流程獎勵建模（Process Reward Modeling, PRM）等技術，建立由規格理解、測試案例生成、測試程式產生、實際API執行、工作流程決策至結果分析之完整自動化測試流程。
\text 研究結果顯示，僅依賴隨機輸入或傳統規格導向方法所產生之測試請求，容易因缺乏語意資訊而導致大量無效請求。透過OpenAPI規格解析與Seeded Input Generation機制，可有效提升測試輸入品質，降低不合理參數組合所造成之失敗案例，顯示規格導向語意輸入建構對RESTful API測試具有實質價值。
\text 在測試案例生成階段，本研究導入多代理協作機制，透過品質代理、開發代理與協調代理共同參與規格分析與約束推理，使測試案例生成不再僅依賴單一大型語言模型之推論結果。實驗結果顯示，多代理討論機制能夠協助系統發現規格中的潛在依賴關係與隱含限制條件，進而產生較符合實際業務邏輯之測試案例與請求序列。此結果驗證多代理協作於API測試領域中不僅能作為內容生成工具，更能作為語意知識整合與規格推理機制。
\text 此外，本研究將工作流程決策問題形式化為馬可夫決策過程（Markov Decision Process, MDP），並建立工作流程決策策略模組，使系統能根據HTTP狀態碼、錯誤類型、依賴狀態與資源可用性等執行訊號，自動選擇適當之恢復行動。研究結果顯示，工作流程決策與恢復機制的主要效益體現在執行後工作流程延續能力上。於100筆主實驗中，v5之Strict Pass Rate與v3、v4同為0.69，但Final Pass Rate提升至0.79，新增之10筆通過全部來自Workflow Recovered Pass，且皆集中於404資源依賴缺失情境。此結果表示，本研究目前最直接被驗證之核心貢獻，在於後執行工作流程恢復層能避免測試流程因單一錯誤而提前終止，進一步提高整體測試流程之可完成度。
\text 在離線強化學習部分，本研究透過工作流程轉移事件建立Offline Reinforcement Learning Dataset，並利用表格式擬合Q值迭代（Tabular Fitted Q Iteration）訓練Offline-Q Policy。實驗結果顯示，Offline-Q Policy已能成功學習部分工作流程決策模式，尤其是在成功後優先進入報告，以及在404資源依賴缺失情境下優先選擇資源解析等局部模式。然而，由於訓練資料規模、行動覆蓋率與狀態空間覆蓋度仍有限，其整體表現尚不足以作為全面優於Rule-based Policy之證據。儘管如此，本研究已驗證工作流程決策問題可被轉換為離線強化學習問題，並證明Workflow Decision Policy具備進一步導入學習型決策策略之可行性。

\text 整體而言，本研究成功建立一套可實際執行之多代理RESTful API自動化測試框架，並驗證大型語言模型、多代理協作機制、後執行工作流程恢復層以及工作流程決策形式化於API測試自動化領域之整合可行性，為未來智慧化API測試系統之發展提供具體參考架構。

\section{研究貢獻與限制}

\subsection{研究貢獻}
\text 本研究之第一項貢獻在於建立後執行工作流程恢復層（Post-Execution Workflow Recovery Layer）。過去多數RESTful API測試研究主要聚焦於單一請求之產生、輸入參數探索或測試覆蓋率提升，而較少關注測試流程在面對實際執行錯誤時之後續恢復行為。本研究將API測試視為一個包含多個狀態轉移與決策行動之動態流程，使系統在404資源缺失等情境下，仍可透過資源解析與再次執行繼續完成任務。
\text 第二項貢獻在於建立結合大型語言模型、多代理系統、MCP與A2A之協作式RESTful API測試架構。相較於傳統單一模型生成測試案例之方法，本研究透過多代理協作機制將規格分析、依賴推理、測試審查與決策分析等任務進行職責分離，使系統能以較結構化方式處理複雜規格理解與測試案例生成問題，同時提高測試流程之可追蹤性與可維護性。
\text 第三項貢獻在於將工作流程決策問題形式化為馬可夫決策過程。透過將API執行後之狀態、候選行動、狀態轉移與獎勵納入統一表示，本研究使工作流程決策不再只是經驗性錯誤處理，而能以可重播、可比較且可學習之序列決策問題加以分析。
\text 第四項貢獻在於驗證學習型工作流程決策策略之可行性。雖然目前Offline-Q策略尚未構成全面優於Rule-based Policy之證據，但研究結果顯示，工作流程決策問題可被轉換為離線強化學習問題，且學習型策略已能在既有資料覆蓋下學得部分局部有效之恢復模式。
\text 第五項貢獻在於提出適用於工作流程決策問題之Process Reward Modeling設計。不同於僅評估最終測試結果之獎勵設計方式，本研究將每一次狀態轉移視為評估單位，透過狀態改善程度、錯誤可處理性與語意資訊增益等指標建立轉移獎勵機制，使系統能評估中間決策步驟之品質。此方法提供Workflow Decision Learning一種較符合實際測試流程之評估方式。
\text 此外，本研究亦得到一項重要觀察：真實Live Execution環境下所產生之執行結果具有高度隨機性與環境依賴性，因此較適合作為描述性能力展示與失敗型態觀察來源，而不宜直接作為不同演算法效能優劣之強證據。另一方面，實際執行過程所蒐集之失敗案例集合能夠反映真實系統之錯誤分布，而經整理與標準化後之測試集合則具備可重播與可比較之特性。此發現對未來RESTful API測試Benchmark之建構具有重要參考價值。

\subsection{研究限制}
\text 為使本研究之限制分析更具系統性，本文將可能影響結果解讀與跨時間重現性之干擾因素整理如下。首先，Live API本身具有狀態變動特性，不同時間點之回應內容、資源可用性與服務端行為可能發生差異。其次，授權資訊可能因Token過期而失效，進而影響部分需驗證身分之操作。再次，外部服務不穩定與API rate limit可能造成暫時性失敗，使同一測試案例於不同執行輪次中產生不同結果。此外，前次可用之測試資料或資源識別碼亦可能於後續實驗中被刪除或失效。另一方面，不同版本在測試程式實體化程度上的差異，亦可能影響表面通過率與執行嚴格度。最後，若OpenAPI規格與實際服務行為不一致，則可能產生規格漂移問題。上述因素雖不必然改變研究方法本身之有效性，但會影響Live API情境下之結果穩定度與跨時間重現性。
\text 雖然本研究驗證所提出方法之可行性，但仍存在若干限制。首先，本研究主要採用TMDB資料集作為實驗環境。雖然TMDB提供完整OpenAPI規格與真實服務環境，且具備足夠之端點數量與跨操作依賴關係，但其本質仍屬於媒體內容查詢領域之應用服務。因此，本研究所得結果未必能完全代表金融、醫療、工業控制或企業內部系統等其他領域API之行為特性。
\text 其次，本研究主要驗證單一API領域環境下之工作流程決策能力。不同API服務在認證方式、業務規則、資料結構與錯誤模式上可能存在顯著差異，因此目前尚無法直接證明所提出之Workflow Decision Policy具備跨領域泛化能力。未來仍需於更多公開API平台與企業系統環境中進一步驗證。
\text 第三，本研究所建立之Offline Reinforcement Learning Dataset規模仍相對有限。由於資料來源主要來自工作流程執行過程中所蒐集之Workflow Decision Events，因此狀態空間與行動空間之覆蓋率仍不足以支撐更複雜之強化學習模型訓練。部分狀態與行動組合於訓練資料中出現次數有限，導致離線策略學習結果仍呈現較高程度之稀疏性。
\text 此外，本研究觀察到Rule-based Policy於工作流程決策問題中已形成相當強勢之基準模型。由於許多RESTful API錯誤模式具有明顯且可解釋之對應關係，例如401錯誤通常對應授權更新、404錯誤通常對應資源解析等，因此基於專家知識建立之規則式策略已能達成相當穩定之決策效果。另一方面，本輪主實驗中所有成功恢復案例皆集中於404資源依賴缺失情境，表示目前被驗證之恢復能力仍偏向特定錯誤型態，而非全面覆蓋所有恢復型行動。此現象使得Offline-Q策略即使成功學習部分決策模式，也未必能於有限資料規模下取得顯著優勢。因此，較穩健之解讀應為：本研究目前已證明學習型策略可於特定恢復情境中提供額外效益，但尚不足以宣稱其已全面超越Rule-based baseline。
\text 最後，由於本研究採用Live API執行環境，結果仍可能受授權狀態、資源可用性、外部服務波動與規格漂移影響。即使同一測試案例於不同時間點執行，也可能產生不同回應內容與恢復結果。因此，Live API情境雖有助於反映真實環境下之測試挑戰，但其結果亦較不易完全重現，未來仍需配合更完整的Benchmark與可重播資料集加以補強。

\section{未來研究方向}
\text 未來研究可朝向更具代表性與可重播性之Benchmark建構方向發展。研究過程中觀察到，真實API執行環境雖然能反映實際系統運作情況，但其結果容易受到服務端更新、外部環境變動與資料狀態改變之影響。因此，可考慮建立Observed Set與Curated Set雙層資料架構。Observed Set用於持續蒐集真實環境中所觀察到之失敗案例與執行事件，以反映真實錯誤分布；Curated Set則透過整理、標準化與分類後形成可重播之Benchmark資料集，使不同研究能在一致條件下進行比較與驗證。
\text 在強化學習方面，未來可擴大Workflow Decision Dataset之規模與多樣性。隨著狀態空間覆蓋率提高，可逐步由目前之表格式離線學習方法擴展至函數逼近（Function Approximation）或深度強化學習（Deep Reinforcement Learning）架構，以提升策略對高維度狀態空間之處理能力。同時亦可探索離線策略評估（Offline Policy Evaluation）與保守型離線強化學習方法，以降低策略學習過程中之分布偏移問題。

\paragraph{成本感知漸進式資料集擴充}
\text 針對當前離線資料集動作覆蓋不足之限制，未來研究可引入成本感知之漸進式資料集擴充機制（incremental dataset expansion），以在控制測試成本的前提下逐步拓寬動作支撐集。其核心思路並非在單次實驗中解鎖完整動作空間，而是設計基於不確定性的探索機制（uncertainty-aware exploration）。在每次自動化測試執行過程中，系統可根據當前Fitted-Q模型對各狀態之估計信心，識別Q值估計較不穩定之狀態，並於這些關鍵節點選擇性地嘗試非主流行動。新產生之轉移紀錄$(s, a, r, s')$可再納入歷史資料集，供下一輪訓練使用。隨著執行輪次累積，資料集得以漸進式擴大，逐步活化Q-table在低頻行動上的估計能力，同時將對服務端環境之衝擊控制在可接受範圍內。
\text 此方向本質上屬於Batch RL架構下之主動資料蒐集問題（active data collection in offline RL），其實施仍需解決兩個前置問題。其一為不確定性估計方法之選擇。表格式Q學習本身不直接提供預測變異數，因此可能需要引入ensemble方法或bootstrap估計。其二為探索行動之成本上界設定，也就是在何種條件下允許系統嘗試非主流行動，而不致引發服務端副作用。上述問題仍待後續研究結合具體API測試場景加以設計與驗證。

\text 此外，本研究目前採用離線強化學習架構進行策略訓練，未來可進一步研究線上強化學習（Online Reinforcement Learning）機制，使系統能於實際執行過程中持續更新決策策略。透過即時回饋與增量學習機制，系統可根據不同API環境之變化動態調整恢復策略，提高長期適應能力。
\text 另一值得探討之方向為情境式多臂老虎機（Contextual Bandit）模型。由於工作流程決策問題本身具有明顯之單步決策特性，部分恢復行動之效果可於短期內立即觀察，因此Contextual Bandit可能比完整強化學習架構更符合部分Workflow Decision場景。未來可進一步比較Contextual Bandit與Offline Reinforcement Learning於工作流程決策問題中之適用性與成本效益。
\text 在泛化能力方面，未來研究可將實驗環境擴展至多個不同領域之公開API與企業系統，包括電子商務平台、金融服務平台、物聯網管理平台及企業資源規劃系統等。透過跨領域驗證，可進一步評估多代理推理機制、Workflow Decision Policy以及Process Reward Modeling之通用性與可擴展性。
\text 最後，隨著代理系統逐漸朝向標準化與互通化方向發展，未來可進一步研究A2A與MCP於多代理測試環境中的標準化應用模式。目前代理系統仍多以框架內部協定進行整合，不同平台之代理間缺乏一致之互通機制。若能建立以A2A負責代理協作、MCP負責工具與上下文存取之標準化代理生態系統，將有助於未來形成可重用、可擴充且具互操作性之智慧化測試平台，進一步推動大型語言模型與代理技術於軟體測試領域之實務應用與研究發展。

\end{ZhChapter}
